\documentclass[10pt,a4paper]{amsart}
\usepackage[margin=19mm]{geometry}
\usepackage{amsmath,amssymb,mathtools}
\usepackage[scr=rsfs,cal=euler]{mathalfa}
\usepackage{xfrac}
\usepackage[unicode,hidelinks,bookmarks=false]{hyperref}
\newcommand{\ie}{i.e.\ }
\newcommand{\cf}{cf.\ }
\newcommand{\resp}{respectively\ }
\newcommand{\Subsection}{\subsection}
\providecommand{\nobreakdash}{\nobreak\hbox{-}}
\setlength{\emergencystretch}{2em}
\multlinegap0pt
\usepackage{bm}
\usepackage{bbm}
\usepackage{dsfont}
\usepackage{xspace}
\usepackage{cancel}
\usepackage{mathtools}
\usepackage{amsthm}
\makeatletter
\@ifundefined{thm}{\newtheorem{thm}{Thm}}{}
\@ifundefined{defn}{\newtheorem{defn}[thm]{Definition}}{}
\@ifundefined{prop}{\newtheorem{prop}[thm]{Proposition}}{}
\makeatother
\title[$\hat{H}$-eigenvalues of Hermitian tensors and some applicat]{$\hat{H}$-eigenvalues of Hermitian tensors and some applications}
\author{Haojie Chen, Yang Yang}
\date{Source: arXiv:2508.12476v1 (CC BY 4.0)}
\begin{document}
\maketitle

\noindent\emph{Source and licence.} $\hat{H}$-eigenvalues of Hermitian tensors and some applications. Version arXiv:2508.12476v1 (CC BY 4.0), distributed under the Creative Commons Attribution 4.0 International licence, as recorded in the arXiv licence metadata for this exact version and shown on the abstract page https://arxiv.org/abs/2508.12476v1. Full rights notice: licenses/arxiv-2508.12476-CC-BY-4.0.txt.\par\smallskip

\noindent\textit{Editorial scope.} Counted unit: the Prop whose source label is \texttt{None} in the article (source0.tex lines 369--371, source label \texttt{None}), delivered together with the article's own complete original proof of it (source0.tex lines 372--407, one \texttt{proof} environment of 36 lines). No mathematics is written, repaired or re-derived: statement and proof are the article's own text line for line. Required context retained uncounted: 4.13 (lines 358--362). Every definition, display and label of those lines is retained unchanged. Omitted from this package and not redistributed: 1--357, 363--368, 408--552. Nothing inside a retained excerpt is omitted or altered: every retained body is delivered exactly as the article's lines are written, so no omission receipt of any kind accompanies this package. Citations: the retained excerpts carry no citation command at all, so no bibliography entry of the article is transcribed and no reference is redistributed. Reference adaptation: none was needed and none was made; every reference inside the delivered unit resolves inside it, so no unresolved reference or citation is produced. Printed slips kept literal: no printed word, symbol, hypothesis or number of the counted statement or of its proof was altered. Layout and class: the article's own class is \texttt{amsart} with packages including amssymb, amsfonts, xy, enumerate, mathrsfs, geometry, amsmath, latexsym, bm, indentfirst, CJK, amsthm, hyperref, cite; the wrapper is the lane's pinned \texttt{amsart} editorial wrapper at 10pt on A4, which restates the theorem-like environments the retained text needs and adds the article's own macro definitions that the retained text needs (none), with extra wrapper packages (bm, bbm, dsfont, xspace, cancel, mathtools). The delivered numbering is the unit's own. Assets: no retained span contains a figure environment, a graphics-inclusion command, a PDF-page inclusion, a table or a TikZ picture; no image file is copied into this package, the package declares no asset dependency and no image file is redistributed with it. Licence: version arXiv:2508.12476v1 of this article is distributed under the Creative Commons Attribution 4.0 International licence, as recorded in the arXiv licence metadata for this exact version and shown on the abstract page https://arxiv.org/abs/2508.12476v1. A full rights notice is delivered at licenses/arxiv-2508.12476-CC-BY-4.0.txt. Characters: every retained excerpt is pure ASCII, so the delivered engine, XeLaTeX with its default Latin Modern text font, reports no missing character.
\begin{defn}
$A$ is called an LL tensor if al the diagonal entries $a_{i\cdots i\bar{i}\cdots\bar{i}}\geq0,$ and for $1\leq i,j\leq n, i\neq j,$
	\begin{align} (a_{i\cdots i\bar{i}\cdots\bar{i}}-\hat{r}_{i}(A))a_{j\cdots j\bar{j}\cdots\bar{j}}\geq r'_{i}(A)r_{j}(A).\label{4.13} \end{align}
	If furthermore, $a_{i\cdots i\bar{i}\cdots\bar{i}}>0$ for $1\leq i\leq n$ and strict inequality holds in the above inequality, then A is called a strict LL tensor.
\end{defn}

\begin{prop}
Let $A$ be a $2m$-th order $n$-dimensional Hermitian tensor. If A is (strictly) diagonally dominated, then A is an (strict) LLK tensor. If A is an (strict) LLK tensor, then A is an (strict) LL tensor.
	\end{prop}

\begin{proof}
Suppose that A is a diagonally dominated tensor. Then for $1\leq i\leq n, a_{i\cdots i\bar{i}\cdots\bar{i}}\geq 0.$ Furthermore, for $i\neq j,$  we have
\begin{align} a_{j\cdots j\bar{j}\cdots\bar{j}}\geq r_{j}(A) \label{4.10}\end{align} and
$a_{i\cdots i\bar{i}\cdots\bar{i}}\geq r_{i}(A)=r_{i}^{j}(A)+|a_{ij\cdots j\bar{j}\cdots \bar{j}}|,$
i.e., \begin{align}  a_{i\cdots i\bar{i}\cdots\bar{i}}-r_{i}^{j}(A)\geq |a_{ij\cdots j\bar{j}\cdots \bar{j}}|.\label{4.11}\end{align}
Multiplying (\ref{4.10}) and (\ref{4.11}), we directly get that $A$ is an LLK tensor.

Next, suppose that A is an LLK tensor. Then either A is diagonally dominated or A is not diagonally dominated.
If A is diagonally dominated, for all $1\leq i, j\leq n$, we have
$a_{j\cdots j\bar{j}\cdots\bar{j}}\geq r_{j}(A)$ and
$a_{i\cdots i\bar{i}\cdots\bar{i}}-\hat{r}_{i}(A)\geq r'_{i}(A).$
Multiplying them we get that $$(a_{i\cdots i\bar{i}\cdots\bar{i}}-\hat{r}_{i}(A))a_{j\cdots j\bar{j}\cdots\bar{j}}\geq r'_{i}(A)r_{j}(A).$$
So $A$ is an LL tensor.

Assume that A is not a diagonally dominated tensor but an LLK tensor. Then there is a $1\leq i_{0}\leq n$ such that
$0\leq a_{i_{0}\cdots i_{0}\bar{i}_0\cdots \bar{i}_0}<r_{i_{0}}(A).$
As $$(a_{i\cdots i\bar{i}\cdots\bar{i}}-r_{i}^{j}(A))a_{j\cdots j\bar{j}\cdots\bar{j}}\geq r_{j}(A)|a_{ij\cdots j\bar{j}\cdots \bar{j}}|$$ for all $i\neq j$ by definition, we have
$a_{i\cdots i\bar{i}\cdots\bar{i}}-r_{i}^{j}(A)\geq |a_{ij\cdots j\bar{j}\cdots \bar{j}}|$, for $i\neq i_0$. So $$a_{i\cdots i\bar{i}\cdots\bar{i}}\geq r_{i}^{j}(A)+|a_{ij\cdots j\bar{j}\cdots \bar{j}}|=r_{i}(A)=\hat{r}_{i}(A)+r'_{i}(A).$$ 
Thus, we get that $(a_{i\cdots i\bar{i}\cdots\bar{i}}-\hat{r}_{i}(A))a_{j\cdots jj}\geq r'_{i}(A)r_{j}(A)$ for all $i,j\neq i_{0}$ for $i, j\neq i_0$. So we only need to deal with the pairs $(i_0, j)$ and $(i, i_0)$ for $i, j\neq i_0$.
For $j\neq i_{0}$, we have
$$(a_{i_{0}\cdots i_{0}\bar{i}_0\cdots \bar{i}_0}-r_{i_{0}}^{j}(A))a_{j\cdots j\bar{j}\cdots \bar{j}}\geq r_{j}(A)|a_{i_{0}j\cdots j\bar{j}\cdots \bar{j}}|. $$
As $0\leq a_{i_{0}\cdots i_{0}\bar{i}_0\cdots \bar{i}_0}<r_{i_{0}}(A),$ we get $a_{i_{0}\cdots i_{0}\bar{i}_0\cdots \bar{i}_0}-r_{i_{0}}^{j}(A)<|a_{i_{0}j\cdots j\bar{j}\cdots \bar{j}}|, a_{i_{0}\cdots i_{0}\bar{i}_0\cdots \bar{i}_0}-\hat{r}_{i_{0}}(A)<r'_{i_{0}}(A).$ If $r_{j}(A)>0,r'_{i_{0}}(A)>0$ then
$$\dfrac{a_{i_{0}\cdots i_{0}\bar{i}_0\cdots \bar{i}_0}-\hat{r}_{i_{0}}(A)}{r'_{i_0}(A)}\dfrac{a_{j\cdots j\bar{j}\cdots \bar{j}}}{r_{j}(A)}\geq \dfrac{a_{i_{0}\cdots i_{0}\bar{i}_0\cdots \bar{i}_0}-r_{i_{0}}^{j}(A)}{|a_{i_{0}j\cdots j\bar{j}\cdots \bar{j}}|}\dfrac{a_{j\cdots j\bar{j}\cdots \bar{j}}}{r_{j}(A)}\geq 1,$$
which implies that (\ref{4.13}) holds for$(i_0, j)$. If $r_{j}(A)=0$ or $r'_{i_{0}}=0$, we have
$$(a_{i_{0}\cdots i_{0}\bar{i}_0\cdots \bar{i}_0}-\hat{r}_{i_{0}(A)})a_{j\cdots j\bar{j}\cdots \bar{j}}\geq 0=r'_{i_{0}}(A)r_{j}(A).$$
So (\ref{4.13}) holds for $i=i_{0}$ and j.
Assume that $i\neq i_{0}.$ By definition we have 
$$(a_{i\cdots i\bar{i}\cdots \bar{i}}-r_{i}^{i_0}(A))a_{i_0\cdots i_0\bar{i}_0\cdots\bar{i}_0}\geq r_{i_0}(A)|a_{ii_0\cdots i_0\bar{i}_0\cdots\bar{i}_0}|.$$
Then $a_{i\cdots i\bar{i}\cdots\bar{i}}\geq r_{i}(A)$ for $i\neq i_{0}$ which gives that $a_{i\cdots i\bar{i}\cdots\bar{i}}-\hat{r}_{i}(A)\geq r'_{i}(A)\geq0.$ If $r_{i_0}(A)>0,r'_{i}(A)>0$, then
$$\dfrac{a_{i\cdots i\bar{i}\cdots\bar{i}}-\hat{r}_{i}(A)}{r'_{i}(A)}\dfrac{a_{i_{0}\cdots i_{0}\bar{i}_0\cdots \bar{i}_0}}{r_{i_{0}}(A)}\geq \frac{a_{i\cdots i\bar{i}\cdots \bar{i}}-r_{i}^{i_0}(A)}{|a_{ii_0\cdots i_0\bar{i}_0\cdots\bar{i}_0}|}\frac{a_{i_0\cdots i_0\bar{i}_0\cdots \bar{i}_0}}{r_{i_0}(A)}\geq 1$$
which implies that (\ref{4.13}) holds for $(i, i_{0}).$ If $r_{i_0}(A)=0$ or $r'_i(A)=0$, then 
$$(a_{i\cdots i\bar{i}\cdots\bar{i}}-\hat{r_{i}}(A))a_{i_{0}\cdots i_{0}\bar{i}_0\cdots \bar{i}_0}\geq0=r'_{i}(A)r_{i_{0}}(A).$$
So(\ref{4.13}) still holds for $(i, i_{0})$. 
Thus, (\ref{4.13}) always holds for $i\neq j$ and $A$ is an LL tensor.
The proofs for strict tensors are similar. 
\end{proof}	

\end{document}
